\documentclass[10pt,a4paper]{amsart}
\usepackage[margin=19mm]{geometry}
\usepackage{amsmath,amssymb,mathtools}
\usepackage[scr=rsfs,cal=euler]{mathalfa}
\usepackage{xfrac}
\usepackage[unicode,hidelinks,bookmarks=false]{hyperref}
\newcommand{\ie}{i.e.\ }
\newcommand{\cf}{cf.\ }
\newcommand{\resp}{respectively\ }
\newcommand{\Subsection}{\subsection}
\providecommand{\nobreakdash}{\nobreak\hbox{-}}
\setlength{\emergencystretch}{2em}
\multlinegap0pt
\usepackage[shortlabels]{enumitem}
\usepackage{amssymb}
\usepackage{xcolor}
\newtheorem{lemma}{Lemma}
\newtheorem{proposition}{Proposition}
\newtheorem{theorem}{Theorem}
\title{Almost Golomb Sequences}
\author{Benoit Cloitre}
\date{Source: arXiv:2604.02404v1 (CC BY 4.0)}
\begin{document}
\maketitle

\noindent\emph{Source and licence.} Almost Golomb Sequences. Version arXiv:2604.02404v1 (CC BY 4.0), distributed by arXiv under the Creative Commons Attribution 4.0 International licence, http://creativecommons.org/licenses/by/4.0/, as recorded in the arXiv licence metadata for this exact version and shown on the abstract page https://arxiv.org/abs/2604.02404v1. Full licence notice: \texttt{licenses/arxiv-2604.02404-CC-BY-4.0.txt}.\par\smallskip

\noindent\textit{Editorial scope.} Counted unit: the article's theorem thm:r3 (source0.tex lines 944--951). The article's complete original proof of it is delivered with it (source0.tex lines 953--1117, one \texttt{proof} environment of 165 lines). No mathematics is written, repaired or re-derived: the statement and the proof are the article's own lines, in the article's own notation, and no retained line is substituted or re-flowed.  Retained uncounted as the source-local context the proof needs: lines 356--361; lines 370--381; lines 394--399; lines 784--793; lines 787--791; lines 839--841; lines 934--936.  Nothing was omitted from any retained span, so the package carries no omission record.  No retained line contains a citation, so the wrapper carries no bibliography and no bibliography entry.  No retained line contains a figure environment, an image-inclusion command or drawing code, so no image is delivered and the package declares no asset dependency.  Editorial formatting only, in addition: an AMS \texttt{amsart} preamble that restates the article's own declarations and macros used by the retained lines, namely the environments \texttt{lemma}, \texttt{proposition}, \texttt{theorem} on separate counters, together with the standard packages those lines need.  The retained environments print in this document as Lemma 1, Proposition 1, Theorem 1 in order of appearance, whereas the article prints the same items with the numbering of its own full document.  The complete native source of the article as distributed in the arXiv e-print for this exact version is bundled unchanged as \texttt{sources/197734/source0.tex}.
\begin{lemma}\label{lem:nested-anchor}
For every $r\ge 2$ and $n\ge 1$,
\[
  S_{S_n}=rn-R_n, \qquad R_n:=\sum_{j=1}^{r-1}(r-j)\,d(S_n-j)\in[0,\,r(r-1)/2].
\]
\end{lemma}

\begin{lemma}\label{lem:local_projection}
Let $A\ge 3$.
\begin{enumerate}[label=(\roman*)]
  \item For $X\ge 0$: $d(S_A+X)=1\iff \sum_{j=0}^{\delta-1}L_{A+j}=X+1$
    for some $\delta\in\{1,\ldots,X+1\}$.
    Hence $d(S_A+X)$ is determined by $(L_A,\ldots,L_{A+X})$.
  \item For $Y\ge 1$: $d(S_A-1-Y)=1\iff \sum_{j=1}^{\delta}L_{A-j}=Y$
    for some $\delta\in\{1,\ldots,Y\}$.
    Hence $d(S_A-1-Y)$ is determined by $(L_{A-1},\ldots,L_{A-Y})$.
    Also, $d(S_A-1)=1$ always.
\end{enumerate}
\end{lemma}

\begin{lemma}\label{lem:anchor-characterization}
For every $r\ge 2$, $m\ge 3$, and every integer $K\ge 1$,
\[
  a(K)=m \iff S_m \le K < S_{m+1}.
\]
\end{lemma}

\begin{proposition}\label{prop:eps-auto}
The sequence $\varepsilon$ satisfies $\varepsilon(n)=0$ for $n<5$, $\varepsilon(5)=1$,
and for every $m\ge 2$:
\begin{equation}\label{eq:eps-rec}
  \varepsilon(3m)=\varepsilon(m-1),\qquad
  \varepsilon(3m+1)=\varepsilon(m),\qquad
  \varepsilon(3m+2)=\varepsilon(m).
\end{equation}
In particular, $\varepsilon$ is $3$-automatic.
\end{proposition}

\begin{equation}
  \varepsilon(3m)=\varepsilon(m-1),\qquad
  \varepsilon(3m+1)=\varepsilon(m),\qquad
  \varepsilon(3m+2)=\varepsilon(m).
\end{equation}

\begin{lemma}\label{lem:r3-gap}
For every $k\ge 2$, $\varepsilon(n)=0$ for all $3^k-1\le n\le 5\cdot 3^{k-1}$.
\end{lemma}

\begin{lemma}\label{lem:subdiag}
For every $r\ge 2$ and every $n\ge 3$, $a(n)\le n-1$.
\end{lemma}

\begin{theorem}\label{thm:r3}
For every $n\ge 2$:
\begin{align}
  a(3n)   &= a(n-2)+a(n-1)+a(n)+1+\varepsilon(n-1), \label{eq:r3a}\\
  a(3n+1) &= a(n-1)+a(n)+a(n+1),                    \label{eq:r3b}\\
  a(3n+2) &= a(n)+a(n+1)+a(n+2)-1-\varepsilon(n).   \label{eq:r3c}
\end{align}
\end{theorem}

\begin{proof}
We prove the three identities \eqref{eq:r3a}--\eqref{eq:r3c} together with
the \emph{run-length invariant}: defining
\begin{equation}\label{eq:Ldef}
  L_n\;:=\;S_{n+1}-S_n\;=\;a(n+1)-a(n-2)\;=\;d(n-2)+d(n-1)+d(n)
  \;\in\;\{1,2,3\}
\end{equation}
(the length of the run of value~$n$, for $n\ge 3$), we prove that for every $n\ge 3$,
\begin{equation}\label{eq:r3-Linv}
  L_n \le 2+\varepsilon(n-1),
  \qquad\text{with equality } L_n=3 \text{ whenever } \varepsilon(n-1)=1.
\end{equation}
The proof is by strong induction on~$n$.
The hypothesis $\mathcal{H}_n$ asserts that the three denesting formulas hold
at all ranks $m<n$, and that the invariant~\eqref{eq:r3-Linv} holds for
all $m$ with $3\le m\le n+1$.
(The invariant at $m\le n+1$ involves only $a(k)$ for $k\le m+1\le n+2<3n$,
hence values already determined at a previous induction step.)

\textit{Base cases.}
Direct computation for $2\le n\le 30$ verifies the denesting formulas
and the invariant~\eqref{eq:r3-Linv}.
This threshold covers all $n$ up to well past $I_1=[16,18]$,
the first non-trivial component of $\mathcal{I}$. For $n\ge 31$, every index
encountered in the induction step satisfies the hypotheses of
Proposition~\ref{prop:eps-auto} (valid for $m\ge 2$) and
Lemma~\ref{lem:r3-gap} (valid for $k\ge 2$).

\textit{Induction step.}
Fix $n\ge 31$ and assume $\mathcal{H}_n$.
We determine the three new values $a(3n)$, $a(3n+1)$, $a(3n+2)$
via Lemma~\ref{lem:anchor-characterization}: for $M\ge 3$,
\begin{equation}\label{eq:anchor-char}
  a(K)=M \iff S_M\le K<S_{M+1}.
\end{equation}
For each candidate $M$ below, one has $M<3n$.
Indeed, by Lemma~\ref{lem:subdiag} ($a(t)\le t-1$ for $t\ge 3$):
\[
  M_1 = S_{n+1} \le n+(n{-}1)+(n{-}2)=3n{-}3,
\]
\[
  M_0 = S_n+1+\varepsilon(n{-}1) \le (n{-}1)+(n{-}2)+(n{-}3)+2 = 3n{-}4,
\]
\[
  M_2 = S_{n+2}-1-\varepsilon(n) \le (n{+}1)+n+(n{-}1)-1 = 3n{-}1.
\]
Hence the values $a(M),a(M-1),a(M-2)$ entering $S_M$ are
already determined by the induction hypothesis.

\textit{Case $a(3n+1)$.}
The candidate is $M_1:=a(n-1)+a(n)+a(n+1)=S_{n+1}$.
By the run structure, $a(S_{n+1})=n+1$ and $a(S_{n+1}-1)=n$.
If $L_n\ge 2$ then $a(S_{n+1}-2)=n$; if $L_n=1$ then
$S_{n+1}-2=S_n-1$ and $a(S_n-1)=n-1$.
Hence $S_{M_1}\in\{3n,\,3n+1\}$, so $S_{M_1}\le 3n+1$.
For the upper bound,
$L_{M_1}=a(S_{n+1}+1)-a(S_{n+1}-2)$.
If $L_{n+1}\ge 2$ then $a(S_{n+1}+1)=n+1$;
if $L_{n+1}=1$ then $a(S_{n+1}+1)=n+2$.
Checking all combinations: $S_{M_1+1}=S_{M_1}+L_{M_1}\ge 3n+2>3n+1$.
By~\eqref{eq:anchor-char}, $a(3n+1)=S_{n+1}$,
which is~\eqref{eq:r3b}.

\textit{Case $a(3n)$.}
The candidate is $M_0:=a(n-2)+a(n-1)+a(n)+1+\varepsilon(n-1)=S_n+1+\varepsilon(n-1)$.

\textit{Subcase $\varepsilon(n-1)=0$.}
Then $M_0=S_n+1$ and $L_n\le 2$ by~\eqref{eq:r3-Linv}.
Since $a(S_n)=n$ and $a(S_n-1)=n-1$:
if $L_n=2$ then $a(M_0)=n$, giving
$S_{M_0}=n+n+(n-1)=3n-1$;
if $L_n=1$ then $M_0=S_{n+1}$, $a(M_0)=n+1$, giving
$S_{M_0}=(n+1)+n+(n-1)=3n$.
In both cases $S_{M_0}\le 3n$.
For the upper bound: when $L_n=2$,
$a(M_0+1)=a(S_{n+1})=n+1$, so $L_{M_0}=(n+1)-(n-1)=2$
and $S_{M_0+1}=3n+1$;
when $L_n=1$,
$L_{M_0}=a(S_{n+1}+1)-a(S_n-1)\ge(n+1)-(n-1)=2$
and $S_{M_0+1}\ge 3n+2$.
In both cases $S_{M_0+1}>3n$.

\textit{Subcase $\varepsilon(n-1)=1$.}
Then $M_0=S_n+2$ and $L_n=3$ by~\eqref{eq:r3-Linv}.
All three positions $S_n,\,S_n+1,\,S_n+2$ carry value~$n$, so
$S_{M_0}=n+n+n=3n$.
Moreover $M_0+1=S_n+3=S_{n+1}$, hence $a(M_0+1)=n+1$,
giving $L_{M_0}=(n+1)-n=1$ and $S_{M_0+1}=3n+1$.

In all subcases $S_{M_0}\le 3n<S_{M_0+1}$,
proving $a(3n)=M_0$, which is~\eqref{eq:r3a}.

\textit{Case $a(3n+2)$.}
The candidate is $M_2:=a(n)+a(n+1)+a(n+2)-1-\varepsilon(n)=S_{n+2}-1-\varepsilon(n)$.

\textit{Subcase $\varepsilon(n)=0$.}
Then $M_2=S_{n+2}-1$ and $L_{n+1}\le 2$.
Here $a(M_2)=a(S_{n+2}-1)=n+1$ (last position of the run of $n+1$).
If $L_{n+1}=2$: $a(M_2-1)=n+1$, $a(M_2-2)=a(S_{n+1}-1)=n$,
giving $S_{M_2}=3n+2$.
If $L_{n+1}=1$: $a(M_2-1)=a(S_{n+1}-1)=n$.
\quad If $L_n\ge 2$: $a(M_2-2)=n$, giving $S_{M_2}=3n+1$.
\quad If $L_n=1$: $a(M_2-2)=a(S_n-1)=n-1$, giving $S_{M_2}=3n$.
In all cases $S_{M_2}\le 3n+2$.
For the upper bound, $M_2+1=S_{n+2}$, so
$S_{M_2+1}=S_{S_{n+2}}$.
By Lemma~\ref{lem:nested-anchor},
$S_{S_{n+2}}=3(n+2)-2d(S_{n+2}-1)-d(S_{n+2}-2)=3n+4-d(S_{n+2}-2)$
(using $d(S_{n+2}-1)=1$, Lemma~\ref{lem:local_projection}$(ii)$).
By Lemma~\ref{lem:local_projection}$(ii)$, $d(S_{n+2}-2)=1$
iff $L_{n+1}=1$.
Hence $S_{M_2+1}\ge 3n+3>3n+2$.

\textit{Subcase $\varepsilon(n)=1$.}
Then $M_2=S_{n+2}-2$ and $L_{n+1}=3$ by~\eqref{eq:r3-Linv}.
The run of $n+1$ occupies $[S_{n+1},\,S_{n+1}+2]$, so
$a(M_2)=a(S_{n+1}+1)=n+1$,
$a(M_2-1)=a(S_{n+1})=n+1$,
$a(M_2-2)=a(S_{n+1}-1)=n$,
giving $S_{M_2}=(n+1)+(n+1)+n=3n+2$.
For the upper bound,
$L_{M_2}=a(S_{n+2}-1)-a(S_{n+1}-2)=(n+1)-a(S_{n+1}-2)$.
If $L_n\ge 2$: $a(S_{n+1}-2)=n$, $L_{M_2}=1$;
if $L_n=1$: $a(S_{n+1}-2)=a(S_n-1)=n-1$, $L_{M_2}=2$.
In both cases $S_{M_2+1}=S_{M_2}+L_{M_2}\ge 3n+3>3n+2$.

In all subcases $S_{M_2}\le 3n+2<S_{M_2+1}$,
proving $a(3n+2)=M_2$, which is~\eqref{eq:r3c}.

\textit{Propagation of the run-length invariant.}
With the denesting established at rank~$n$, a direct calculation
using \eqref{eq:r3a}--\eqref{eq:r3c} at ranks $n-1$ and $n$ gives:
\begin{align}
  L_{3n-1} &= L_{n-1}+\varepsilon(n-1)-\varepsilon(n-2),\label{eq:Lprop1}\\
  L_{3n}   &= L_n,                                       \label{eq:Lprop2}\\
  L_{3n+1} &= L_{n+1}+\varepsilon(n-1)-\varepsilon(n).   \label{eq:Lprop3}
\end{align}
The $\varepsilon$-recurrence~\eqref{eq:eps-rec} gives
$\varepsilon(3n-2)=\varepsilon(3n-1)=\varepsilon(3n)=\varepsilon(n-1)$.

\emph{Upper bounds.}
$L_{3n-1}\le(2+\varepsilon(n-2))+\varepsilon(n-1)-\varepsilon(n-2)=2+\varepsilon(n-1)=2+\varepsilon(3n-2)$.
$L_{3n}=L_n\le 2+\varepsilon(n-1)=2+\varepsilon(3n-1)$.
$L_{3n+1}\le(2+\varepsilon(n))+\varepsilon(n-1)-\varepsilon(n)=2+\varepsilon(n-1)=2+\varepsilon(3n)$.

\emph{Exactness when $\varepsilon=1$.}
Suppose $\varepsilon(3n-2)=\varepsilon(n-1)=1$.
By~\eqref{eq:r3-Linv} at~$n$, $L_n=3$, hence $d(n-2)=d(n-1)=d(n)=1$.
Then $L_{n-1}=d(n-3)+d(n-2)+d(n-1)\ge 2$.
Combined with $L_{n-1}\le 2+\varepsilon(n-2)$:
if $\varepsilon(n-2)=1$ then $L_{n-1}=3$ and $L_{3n-1}=3+1-1=3$;
if $\varepsilon(n-2)=0$ then $L_{n-1}=2$ and $L_{3n-1}=2+1-0=3$.

Suppose $\varepsilon(3n-1)=\varepsilon(n-1)=1$.
Then $L_{3n}=L_n=3$ directly.

Suppose $\varepsilon(3n)=\varepsilon(n-1)=1$.
By~\eqref{eq:r3-Linv} at~$n$, $L_n=3$, hence $d(n-1)=d(n)=1$.
Then $L_{n+1}=d(n-1)+d(n)+d(n+1)\ge 2$.
Combined with $L_{n+1}\le 2+\varepsilon(n)$:
if $\varepsilon(n)=1$ then $L_{n+1}=3$ and $L_{3n+1}=3+1-1=3$;
if $\varepsilon(n)=0$ then $L_{n+1}=2$ and $L_{3n+1}=2+1-0=3$.

This closes the strong induction.
\end{proof}

\end{document}
